\section{绪论}

\subsection{研究背景}
低空经济作为国家战略性新兴产业正在快速崛起。据中国民航局预测，到2025年我国低空经济市场规模将达1.5万亿元，无人机注册数量已突破120万架。在物流配送、农业植保、城市巡检、应急救援等领域，无人机的规模化应用正从试点走向常态化运营。然而，低空无人机飞行高度低（通常120--600\,m）、速度慢（10--30\,m/s）、雷达反射截面积小（RCS通常小于0.01\,m$^2$），传统雷达系统难以对其进行有效监管，形成了低空安全的监管盲区。

\subsection{问题提出}
传统的低空监管依赖专用雷达站点，存在部署成本高、覆盖密度低、无法与通信网络联动等问题。以典型的X波段监视雷达为例，单套设备采购及部署费用在数百万元级别，且需要独立的站址、供电和回传链路。而现有5G蜂窝网络在城区的基站间距仅约200\,m，已实现密集覆盖。如何复用这一庞大的通信基础设施实现低空感知功能，成为产学研共同关注的焦点问题。

\subsection{ISAC技术动因}
通信感知一体化（Integrated Sensing and Communication, ISAC）技术为解决上述问题提供了全新思路。其核心思想是在现有5G基站上叠加类似雷达的感知能力：基站在发射通信信号的同时，利用目标反射的回波信号提取距离、速度和角度等感知参数。5G NR采用的OFDM波形天然具备雷达感知能力——其多载波结构支持在子载波维度估计距离、在符号维度估计速度。因此，ISAC可以在不增加额外硬件（或仅需少量改造）的前提下，将通信基站升级为"通信+感知"双功能节点，其技术动因可归纳为三点：

\begin{enumerate}[label=(\arabic*)]
  \item \textbf{覆盖复用}：城区5G基站密度极高（站间距$\sim$200\,m），天然形成密集的感知网络，无需额外建设雷达站。
  \item \textbf{硬件共享}：ISAC复用基站的天线阵列、射频前端和基带处理单元，增量硬件成本极低。
  \item \textbf{数据联动}：感知数据可通过5G核心网即时回传至云端平台，与通信数据、飞行计划数据融合处理。
\end{enumerate}

\subsection{课题目标与报告结构}
本报告沿"低空物联网需求$\to$ISAC感知原理$\to$3GPP标准化$\to$IoT平台架构$\to$产业应用"这一主线，构建完整的技术分析闭环。第二章推导OFDM雷达信号模型与关键算法并进行仿真验证；第三章梳理3GPP标准化进展；第四章设计感知数据接入IoT平台的系统架构；第五章对比分析产业应用案例；第六章讨论挑战与展望。
